\chapter{Results and Discussion}
\label{ch:results}

This chapter presents the results of the experimental evaluation described in Chapter~\ref{ch:experiments} and discusses their implications for the research questions posed in Section~\ref{sec:research-questions}.

\textit{Note: all numeric entries and figure placeholders marked \textbf{[TODO]} are to be filled in after running the system on the 20-URL dataset.}

\section{Clustering Results}
\label{sec:clustering-results}

\subsection{DBCV Scores Across Levels}

The DBCV scores obtained for each level are reported in Table~\ref{tab:dbcv-by-level} (Chapter~\ref{ch:experiments}). Figure~\ref{fig:dbcv-bar} visualises these scores as a bar chart.

\begin{figure}[htbp]
    \centering
    \fbox{\parbox{0.85\textwidth}{\centering\vspace{2.5cm}[TODO: bar chart of DBCV score per level, copied from Chapter 5 / generated by the system]\vspace{2.5cm}}}
    \caption{DBCV score per DOM depth level for the 20-URL dataset.}
    \label{fig:dbcv-bar}
\end{figure}

The automatically selected best level is \textbf{[TODO]}, with a DBCV score of \textbf{[TODO]} and \textbf{[TODO]} clusters. This level was chosen by the median-based algorithm (Section~\ref{sec:similarity-analysis}): valid levels were sorted by cluster count, the cluster count at the median position was used as threshold $W = \text{[TODO]}$, and among levels with at most $W$ clusters the one with the highest DBCV score was returned.

\subsection{Colour-Coded HTML Visualisations}

For each URL the system generates an HTML file in which every DOM element is rendered with the background colour of its assigned cluster. Figure~\ref{fig:html-vis} shows representative screenshots for one page from each category.

\begin{figure}[htbp]
    \centering
    \fbox{\parbox{0.85\textwidth}{\centering\vspace{4cm}[TODO: side-by-side screenshots of colour-coded HTML output for one nvd.nist.gov page, one pirateparty.gr page, and one gastronomos.gr page at the best level]\vspace{4cm}}}
    \caption{Colour-coded HTML visualisations at the best level. Each colour represents one cluster; pages from the same thematic category are expected to share similar colour distributions.}
    \label{fig:html-vis}
\end{figure}

\section{Similarity Matrix Results}
\label{sec:similarity-results}

\subsection{Combined Similarity Scores}

The full inter-domain similarity matrix at the best level is given in Table~\ref{tab:similarity-matrix} (Chapter~\ref{ch:experiments}). Figure~\ref{fig:heatmap-result} renders the same matrix as a Viridis heatmap.

\begin{figure}[htbp]
    \centering
    \fbox{\parbox{0.85\textwidth}{\centering\vspace{4cm}[TODO: heatmap of the 6$\times$6 similarity matrix — screenshot from the React frontend or regenerated with matplotlib]\vspace{4cm}}}
    \caption{Inter-domain similarity heatmap at the best level. Brighter cells indicate higher combined similarity. The security block (nvd / jvn / kb.cert) and the food block (dia-trofis / gastronomos) are expected to appear as bright sub-matrices.}
    \label{fig:heatmap-result}
\end{figure}

\subsection{Discussion of RQ1}

\textit{Fill in after obtaining results. Address: are within-category scores (security--security, food--food) higher than cross-category scores? Provide the observed within-category average and the cross-category average, and state whether the pattern is consistent with the expected groupings. If any unexpected high or low score is observed, refer to the HTML visualisations to explain it.}

\section{Individual Metric Contributions}
\label{sec:individual-metrics}

Beyond the combined score, the three component metrics can be compared to understand their individual contributions.

\begin{table}[htbp]
\centering
\caption{Component similarity scores for selected domain pairs at the best level}
\label{tab:component-metrics}
\begin{tabular}{lcccc}
\toprule
\textbf{Domain pair} & \textbf{Cosine} & \textbf{EMD} & \textbf{Hausdorff} & \textbf{Combined} \\
\midrule
nvd.nist.gov -- jvn.jp         & [TODO] & [TODO] & [TODO] & [TODO] \\
nvd.nist.gov -- kb.cert.org    & [TODO] & [TODO] & [TODO] & [TODO] \\
jvn.jp -- kb.cert.org          & [TODO] & [TODO] & [TODO] & [TODO] \\
dia-trofis.gr -- gastronomos.gr & [TODO] & [TODO] & [TODO] & [TODO] \\
nvd.nist.gov -- pirateparty.gr & [TODO] & [TODO] & [TODO] & [TODO] \\
nvd.nist.gov -- gastronomos.gr & [TODO] & [TODO] & [TODO] & [TODO] \\
\bottomrule
\end{tabular}
\end{table}

\section{Processing Time}
\label{sec:processing-results}

The per-phase processing times are reported in Table~\ref{tab:processing-time} (Chapter~\ref{ch:experiments}). DOM extraction dominates the total time because each URL requires a full headless browser session with computed-style retrieval. Clustering and similarity computation are comparatively fast once the features are available.

\section{Discussion of RQ2 and RQ3}

\textit{Fill in after obtaining results. For RQ2: identify the best level and whether it falls in the mid-range (3--7) as the system's recommended range suggests. For RQ3: compare the DOM extraction time per URL against the clustering and similarity times; note whether the relationship is approximately linear.}

\section{Limitations Observed During Experiments}
\label{sec:limitations-discussion}

\subsection{Dynamic Content}

Pages that load content asynchronously may be captured before all elements are rendered, resulting in incomplete DOM trees. The headless Chrome instance uses a fixed page-load timeout; pages that exceed it are skipped and logged.

\subsection{Single-Browser Rendering}

The system captures computed styles as reported by Chrome. Sites with browser-specific CSS may produce different feature vectors in other browsers. This is a known constraint of any single-browser DOM analysis approach.

\subsection{Dataset Size}

The experimental dataset of 20 URLs is sufficient for a proof-of-concept evaluation across three thematic categories, but is small relative to the scale at which the system is intended to be used. Larger datasets would allow statistical validation of the observed similarity patterns.

\section{Summary}
\label{sec:ch6-summary}

This chapter presented the clustering and similarity results for the 20-URL dataset, including DBCV scores per level, the automatically selected best level, the full inter-domain similarity matrix, and per-phase processing times. The results address the three research questions posed in Chapter~\ref{ch:experiments}: whether the similarity metric reflects thematic groupings (RQ1), which level achieves the best clustering quality (RQ2), and how processing time scales with the dataset (RQ3). Detailed discussion is deferred to the filled-in version of this chapter after running the experiments.
